%
% Standard Issue #20
%   https://github.com/mpi-forum/mpi-issues/issues/20
%
% Formerly Ticket 323
%    https://svn.mpi-forum.org/trac/mpi-forum-web/ticket/323
%
% More information on implementation and development process
%    https://fault-tolerance.org
%
%%%%%%%%%%%%%%%%%%%%%%%%%%%%%%%%%%%%%%%%%%%%%%%%%%%%%%%%%%%%%%%%%%%%%%

\chapter{Process Fault Tolerance}
\label{sec:ft}
\label{chap:ft}
\mpitermtitleindex{error handling!fault tolerance}
\mpitermtitleindex{error handling!process failure}

\section{Introduction}
\label{sec:ft-intro}

\par In distributed systems with numerous or complex components,
a serious risk is that a component fault manifests as a process failure
that disrupts the normal execution of a long running application.
A process failure is a common outcome for many
hardware, network, or software faults that cause a process to crash;
it can be more formally defined as a fail-stop failure: the affected
\MPI/ process unexpectedly and permanently stops communicating.
This chapter introduces \MPI/ features that support the development of
applications, libraries, and programming languages that can tolerate
\MPI/ process failures. The primary goal is to specify error classes and
interfaces that permit users to continue simple \MPI/ communication
(e.g., some point-to-point patterns) after failures have impacted the execution and rebuild
\MPI/ objects (e.g., communicators, etc.) as needed to restore the full
capability of \MPI/ to carry out communication operations (like
collective communications). This specification does not include
mechanisms to restore the application data lost due to process failures. The
literature is rich with diverse fault tolerance techniques that
the users may employ at their discretion, including
checkpoint-restart, algorithmic dataset recovery, and continuation
ignoring failed \MPI/ processes. All these fault tolerance approaches
benefit from, and often require, the definitions and interfaces
specified in this chapter in order to resume communicating after a failure.

\par The expected behavior of \MPI/ in the case of an \MPI/ process failure
is defined by the following statements: any \MPI/ operation that
involves a failed process must eventually either
succeed or raise an \MPI/ error (see
Section~\ref{sec:ft-notification}); an \MPI/ operation that does not
involve a failed process will complete normally, unless interrupted
by the user through provided functionality.
%
If an application
needs comprehensive knowledge of failures, it can use the
interfaces defined in Section~\ref{sec:ft-functions} to explicitly
propagate the notification of locally detected failures, or set
communicators in specific modes that enforce such propagation, see Section~\ref{sec:ft-modes}. %sec:reinit

\par Some usage patterns on reliable machines do not
require fault tolerance. An \MPI/ implementation that does not
tolerate process failures must never raise a \mpitermni{fault tolerance error}\mpitermindex{error handling!fault tolerance!fault tolerance error}
(as listed in Section~\ref{sec:ft-errorcodes}). Applications
using the interfaces defined in this chapter must be portable across
\MPI/ implementations (including those which do not provide
fault tolerance). If an \MPI/ implementation does not provide fault
tolerance features, it may exhibit undefined behavior after a process
failure at any \MPI/ process.
%TODO: a future PR will add a reference to the mechanism used to query
% what FT models are runtime supported by the implementation.

\begin{users}
The \MPI/ standard does not specify transparent process recovery upon \MPI/ process failure.
In particular, restoring the lost dataset, spawning spare processes
or taking other recovery actions are the responsibility of the user.

Many of the operations and semantics described in
this chapter are applicable only when the \MPI/ application has replaced
the default error handler \mpiconst{MPI\_ERRORS\_ARE\_FATAL} on the communicators
it uses.
\end{users}

\section{Failure Notification}
\label{sec:ft-notification}
\mpitermtitleindex{error handling!fault tolerance!notification}

\par All \MPI/ procedures must eventually return (possibly by raising one of the process
failure error classes listed in Section~\ref{sec:ft-errorcodes}) even if they involve a failed \MPI/ process.

\begin{implementors}
As long as an implementation can successfully complete some operations during
a completing procedure, it may choose to delay raising an error. Another valid
implementation might choose to raise an error as quickly as possible.
\end{implementors}

\par When an operation raises a process failure error it may not satisfy its
specification (like any other error, see~\ref{sec:ei-error-classes}).

\subsection{Revoked Communicator}
\mpitermtitleindex{error handling!fault tolerance!revoked}

A communicator is \mpitermdefni{revoked}\mpitermdefindex{error handling!fault tolerance!revoked}
at a given \MPI/ process either when
\mpifunc{MPI\_COMM\_REVOKE} is locally called on it, or when any \MPI/
operation on the communicator raises an error of class
\error{MPI\_ERR\_REVOKED} at that process. Once a
communicator has been revoked at an \MPI/ process, all subsequent non-local operations on
that communicator are considered
local and must complete by raising an error of class
\error{MPI\_ERR\_REVOKED} at that \MPI/ process.

\begin{implementors}
An implementation may detect that a communicator should become revoked
when completing procedures can still complete some operations successfully.
The implementation may delay raising the \error{MPI\_ERR\_REVOKED} error to a later operation,
but it cannot delay revoking the communicator indefinitely and should favor
raising the error promptly.
\end{implementors}


\subsection{Error Reporting Range}
\label{sec:ft-modes}

By default, process failure errors are raised only during communication operations
in which a failed process is involved, but
users can control the range of \MPI/ processes whose failure cause \MPI/ operations
to raise errors on their communicators by setting the following values to the
info key \infokey{mpi\_error\_range} on their communicators:

\begin{description}
\item \infoval{operation}
A process fault tolerance error shall be raised only if the operation involves
a failed \MPI/ process (e.g., a point-to-point message between two non-failed \MPI/ processes
shall not raise a process fault tolerance error).
An \MPI/ process is considered involved in an operation (for the purpose
of this definition) if its failure may prevent it from calling a matching
procedure for the operation, that is, an \MPI/ process is involved in a
communication if any of the following is true:
the process is in the group(s) over which the operation is collective;
the process is a destination or a specified or a matched source in a
point-to-point communication;
or the process belongs to the source group in an \mpiconst{MPI\_ANY\_SOURCE}
receive operation.

\item \infoval{group}
The failure of any \MPI/ process in the group of the communicator causes the communicator to become revoked\mpitermindex{error handling!fault tolerance!revoked}
(as if \mpifunc{MPI\_COMM\_REVOKE} had been called on the communicator).

\item \infoval{universe}
The failure of any \MPI/ process in the \MPI/ universe causes the communicator to become revoked
(as if \mpifunc{MPI\_COMM\_REVOKE} had been called on the communicator).
\end{description}

When the info key \infokey{mpi\_error\_range} is not set, it is equivalent to having the key set to \infoval{operation}.

\subsection{Fault Tolerance Errors in Point-to-Point Communication}
\label{sec:ft-notification:p2p-comm}
\mpitermtitleindex{error handling!fault tolerance!communicator}

An \MPI/ point-to-point communication raises errors of the following
classes to notify users that the communication
could not complete successfully because of the failure of at least one involved \MPI/ process:

\begin{itemize}

\item \error{MPI\_ERR\_REVOKED} indicates that the communicator is revoked.

\item \error{MPI\_ERR\_PROC\_FAILED\_PENDING} indicates, for a nonblocking or persistent communication,
   that the communication is a receive operation from
   \mpiconst{MPI\_ANY\_SOURCE} and no send operation has matched, yet
    a potential sending \MPI/ process has failed. The request remains
    active.

\item In all other cases, the operation raises an error of class
    \error{MPI\_ERR\_PROC\_FAILED} to indicate that the failure prevents the
    operation from following its failure-free specification. If there is a
    request representing a point-to-point communication, it is completed.

\end{itemize}

\subsection{Fault Tolerance Errors in Collective Operations}
\label{sec:ft-notification:coll}

A collective operation raises errors of the following classes to notify
users that the communication could not complete successfully because of
the failure of at least one involved \MPI/ process:

\begin{itemize}

\item \error{MPI\_ERR\_REVOKED} indicates that the communicator is revoked.

\item \error{MPI\_ERR\_PROC\_FAILED} indicates that the failure prevents the
    operation from following its failure-free specification. If there is a
    request representing a collective communication, it is completed.

\end{itemize}


\begin{users}

Depending on how the collective operation is implemented and when an \MPI/ process
failure occurs, some participating \MPI/ processes may raise an error while
other \MPI/ processes return successfully from the same collective operation. For
example, in \mpifunc{MPI\_BCAST}, the root process may succeed before a failed
\MPI/ process disrupts the operation, resulting in some other \MPI/ processes raising an
error.

\end{users}

\begin{users}

Note that communicator creation procedures (e.g., \mpifunc{MPI\_COMM\_DUP} or
\mpifunc{MPI\_COMM\_SPLIT}) are collective operations. As such, if a failure
happened during the procedure, an error might be raised at some \MPI/ processes while
others succeed and obtain a new communicator handle.  Although it is valid to communicate
between \MPI/ processes that succeeded in creating the new communicator handle, the user is
responsible for ensuring a consistent view of the communicator creation, if needed.
\end{users}

\par After an \MPI/ process failure, \mpifunc{MPI\_COMM\_FREE} (as with all
other collective operations) may not complete successfully at all
\MPI/ processes. For any \MPI/ process that receives the return code
\error{MPI\_SUCCESS}, the behavior is defined in
Section~\ref{subsec:context-intracomdest}. If an \MPI/ process raises a process
failure error (classes \error{MPI\_ERR\_PROC\_FAILED} or
\error{MPI\_ERR\_REVOKED}), the communicator handle \mpiarg{comm} is
set to \mpiconst{MPI\_COMM\_NULL}; however, the implementation makes no
guarantee about the success or failure of the
\mpifunc{MPI\_COMM\_FREE} operation, locally or remotely.

\begin{users} Users are encouraged to call \mpifunc{MPI\_COMM\_FREE}
    on communicators they do not wish to use anymore, even when they
    contain failed \MPI/ processes. Although the operation may raise a
    fault tolerance error and not synchronize properly, this
    gives a high quality implementation an opportunity to release
    local resources and memory consumed by the object.
\end{users}

\subsubsection{Error Uniformity}

As noted above, by default, collective communication do not enforce uniformity
in error raising across \MPI/ processes. Despite the performance advantages that
non-uniformity offers, a common usage pattern in applications is
to transform non-uniform error raising into a uniform behavior across all
\MPI/ processes of the group.

Users can set collective operations on a communicator to enforce uniform error
raising by setting the following values in the info key \infokey{mpi\_error\_uniform} on the communicator:

\begin{description}
\item \infoval{local} Process fault tolerance errors are raised to indicate that an \MPI/ process failure prevents from guaranteeing the specified behavior at this process for the collective communication operation (e.g., the output buffer contains invalid data). Other \MPI/ processes may have satisfied their specification (e.g., the output buffer is valid at that process) and may have returned \error{MPI\_SUCCESS}.
\item \infoval{coll} Process fault tolerance errors are raised to indicate that an \MPI/ process failure prevents from guaranteeing the specified behavior at any process for the collective communication operation. Non-synchronizing collective communication become synchronizing.
\item \infoval{construct} Process fault tolerance errors are raised to indicate that an \MPI/ process failure prevents from guaranteeing the specified behavior at any process for the collective \MPI/ communication context constructor/destructor operation (e.g., \mpifunc{MPI\_COMM\_DUP} could not create a new communicator at any process in the group of the communicator). Non-synchronizing context constructor/destructor operations become synchronizing. Collective communication that do not create or free a communication context are not impacted. %TODO: list? what about set_info?
\end{description}

When the info key \infokey{mpi\_error\_uniform} is not set on the communicator, it is equivalent to having the key set to \infoval{local}.
The info key \infokey{mpi\_error\_uniform} must be set to the same value at all \MPI/ processes
for a given communicator.

\subsection{Files and One-sided Operations}
\label{sec:ft-notification:file-rma}

\begin{users}
    Support for fault tolerance with \MPI/ Files and \RMA/ Windows is not
    defined by the \MPI/ standard.
\end{users}

\section{Failure Mitigation Procedures}
\label{sec:ft-functions}
\mpitermtitleindex{error handling!fault tolerance!mitigation}

\subsection{Communicator Procedures}
\label{sec:ft-functions:commfunctions}
\mpitermtitleindex{error handling!fault tolerance!communicator}

\par Process failure notification is not global in \MPI/.
\MPI/ processes that do not call operations involving a failed \MPI/ process are
possibly never notified of its failure (see
Section~\ref{sec:ft-notification}). If a
notification must be propagated, \MPI/
provides a procedure to revoke a communicator at
all members.

\cdeclmainindex{MPI\_Comm}
\begin{funcdef}{MPI\_COMM\_REVOKE}{(\mbox{comm})}
\funcarg{\IN}{comm}{communicator (handle)}
\end{funcdef}
\mpibindmain{MPI\_Comm\_revoke}{(\gb{}MPI\_Comm~comm)}
\mpifnewbindmain{MPI\_Comm\_revoke}{(comm, ierror) \fargs TYPE(MPI\_Comm), INTENT(IN)\ ::\ \mbox{comm}\fnfinalarg{}INTEGER, OPTIONAL, INTENT(OUT)\ ::\ \mbox{ierror}}
\mpifbindmain{MPI\_COMM\_REVOKE}{(COMM, IERROR) \fargs \fnfinalarg{}INTEGER \mbox{COMM}, \mbox{IERROR}}
\mpitermindex{error handling!fault tolerance!revoked}

\par
This procedure notifies all \MPI/ processes in the groups (local and remote)
associated with the communicator \mpiarg{comm} that this communicator
is revoked. The revocation of a communicator by any
\MPI/ process completes non-local \MPI/ operations on
\mpiarg{comm} at all \MPI/ processes by raising an error of class
\error{MPI\_ERR\_REVOKED}.
This procedure is not collective and
therefore does not have a matching call on remote \MPI/ processes. All
non-failed \MPI/ processes belonging to \mpiarg{comm} will be notified of the
revocation despite failures.

\cdeclmainindex{MPI\_Comm}
\begin{funcdef}{MPI\_COMM\_IS\_REVOKED}{(\mbox{comm}, \mbox{flag})}
\funcarg{\IN}{comm}{communicator (handle)}
\funcarg{\OUT}{flag}{\mpicode{true} if the communicator is revoked (logical)}
\end{funcdef}
\mpibindmain{MPI\_Comm\_is\_revoked}{(\gb{}MPI\_Comm~comm, int~*flag)}
\mpifnewbindmain{MPI\_Comm\_is\_revoked}{(comm, flag, ierror) \fargs TYPE(MPI\_Comm), INTENT(IN)\ ::\ \mbox{comm}\fnarg{}LOGICAL, INTENT(OUT)\ ::\ \mbox{flag}\fnfinalarg{}INTEGER, OPTIONAL, INTENT(OUT)\ ::\ \mbox{ierror}}
\mpifbindmain{MPI\_COMM\_IS\_REVOKED}{(COMM, FLAG, IERROR) \fargs INTEGER \mbox{COMM}, \mbox{IERROR}\fnfinalarg{}LOGICAL \mbox{FLAG}}
\mpitermindex{error handling!fault tolerance!revoked}

\par
Returns \mpiarg{flag = true} if the communicator associated with the handle \mpiarg{comm} is revoked at the calling process. It returns \mpiarg{flag = false} otherwise. The operation is local.

\begin{users}
    In a multithreaded application, a thread calling \mpifunc{MPI\_COMM\_IS\_REVOKED} may return \mpiarg{flag = true} before the operation that raises the first exception of class \error{MPI\_ERR\_REVOKED} has completed in a concurrent thread.
\end{users}

\cdeclmainindex{MPI\_Comm}
\cdeclmainindex{MPI\_Group}
\begin{funcdef}{MPI\_COMM\_GET\_FAILED}{(\mbox{comm}, \mbox{failedgrp})}
\funcarg{\IN}{comm}{communicator (handle)}
\funcarg{\OUT}{failedgrp}{group of failed processes (handle)}
\end{funcdef}
\mpibindmain{MPI\_Comm\_get\_failed}{(\gb{}MPI\_Comm~comm, MPI\_Group~*failedgrp)}
\mpifnewbindmain{MPI\_Comm\_get\_failed}{(comm, failedgrp, ierror) \fargs TYPE(MPI\_Comm), INTENT(IN)\ ::\ \mbox{comm}\fnarg{}TYPE(MPI\_Group), INTENT(OUT)\ ::\ \mbox{failedgrp}\fnfinalarg{}INTEGER, OPTIONAL, INTENT(OUT)\ ::\ \mbox{ierror}}
\mpifbindmain{MPI\_COMM\_GET\_FAILED}{(COMM, FAILEDGRP, IERROR) \fargs \fnfinalarg{}INTEGER \mbox{COMM}, \mbox{FAILEDGRP}, \mbox{IERROR}}

\par
This local procedure returns the group \mpiarg{failedgrp} of \MPI/ processes from the
communicator \mpiarg{comm} that are known to have failed by the calling \MPI/ process.
Any \MPI/ process whose failure caused a procedure on \mpiarg{comm} to raise
a process failure error at this \MPI/ process must appear in that group.
The \mpiarg{failedgrp} can be empty, that is, equal to \mpiconst{MPI\_GROUP\_EMPTY}.

When an \MPI/ process is a member of \mpiarg{failedgrp}, it is also a member,
with the same rank, in any group produced by a subsequent call to \mpifunc{MPI\_COMM\_GET\_FAILED} on \mpiarg{comm} at the same calling process.

\begin{users}
\MPI/ makes no assumption about asynchronous
progress of the failure detection. A valid \MPI/ implementation may choose
to update the group of locally known failed \MPI/ processes only when it enters
a procedure that must raise a fault tolerance error.

It is possible that only the calling \MPI/ process has
detected the reported failure.
\end{users}

\cdeclmainindex{MPI\_Comm}
\begin{funcdef}{MPI\_COMM\_ACK\_FAILED}{(\mbox{comm}, \mbox{num\_to\_ack}, \mbox{num\_acked})}
\funcarg{\IN}{comm}{communicator (handle)}
\funcarg{\IN}{num\_to\_ack}{Maximum number of process failures to acknowledge (integer)}
\funcarg{\OUT}{num\_acked}{Number of process failures acknowledged (integer)}
\end{funcdef}
\mpibindmain{MPI\_Comm\_ack\_failed}{(\gb{}MPI\_Comm~comm, int~num\_to\_ack, int~*num\_acked)}
\mpifnewbindmain{MPI\_Comm\_ack\_failed}{(comm, num\_to\_ack, num\_acked, ierror) \fargs TYPE(MPI\_Comm), INTENT(IN)\ ::\ \mbox{comm}\fnarg{}INTEGER, INTENT(IN)\ ::\ \mbox{num\_to\_ack}\fnarg{}INTEGER, INTENT(OUT)\ ::\ \mbox{num\_acked}\fnfinalarg{}INTEGER, OPTIONAL, INTENT(OUT)\ ::\ \mbox{ierror}}
\mpifbindmain{MPI\_COMM\_ACK\_FAILED}{(COMM, NUM\_TO\_ACK, NUM\_ACKED, IERROR) \fargs \fnfinalarg{}INTEGER \mbox{COMM}, \mbox{NUM\_TO\_ACK}, \mbox{NUM\_ACKED}, \mbox{IERROR}}

\par
This local procedure is used to
\mpitermdefni{acknowledge}\mpitermdefindex{error handling!fault tolerance!ack}
locally notified failures on \mpiarg{comm}. The operation acknowledges the first \mpiarg{num\_to\_ack}
process failures on \mpiarg{comm}, that is, it acknowledges the failure of members with a rank lower than
\mpiarg{num\_to\_ack} in the group that would be produced by a concurrent call to \mpifunc{MPI\_COMM\_GET\_FAILED}
on the same \mpiarg{comm}.

The operation also sets the value of \mpiarg{num\_acked} to the current number of acknowledged process
failures in \mpiarg{comm}, that is, a process failure has been acknowledged on \mpiarg{comm} if and only if
the rank of the process is lower than \mpiarg{num\_acked} in the group that would be produced by a subsequent call to
\mpifunc{MPI\_COMM\_GET\_FAILED} on the same \mpiarg{comm}.

\mpiarg{num\_acked} can be larger than \mpiarg{num\_to\_ack} when process failures have been acknowledged in a prior call
to \mpifunc{MPI\_COMM\_ACK\_FAILED}.

After an \MPI/ process failure is acknowledged on \mpiarg{comm}, unmatched \mpiconst{MPI\_ANY\_SOURCE} receive operations
on the same \mpiarg{comm} that would have raised an error of class \error{MPI\_ERR\_PROC\_FAILED\_PENDING}
(see Section~\ref{sec:ft-notification:p2p-comm}) proceed without further
raising errors due to this acknowledged failure.

\begin{users}
    One may query, without side effect, for the number of currently aknowledged
    process failures in \mpiarg{comm} by supplying $0$ in \mpiarg{num\_to\_ack}.
    Conversely, one may unconditionally acknowledge all currently known process
    failures in \mpiarg{comm} by supplying the size of the group of \mpiarg{comm}
    in \mpiarg{num\_to\_ack}.
    Note that the number of acknowledged \MPI/ processes, as returned in \mpiarg{num\_acked},
    can be smaller or larger than the value supplied in \mpiarg{num\_to\_ack}; It is
    however never larger than the size of the group returned by a subsequent call
    to \mpifunc{MPI\_COMM\_GET\_FAILED}.

    Calling \mpifunc{MPI\_COMM\_ACK\_FAILED} on a communicator with failed
    \MPI/ processes has no effect on collective operations.
    If a collective operation would raise an
    error due to the communicator containing a failed process (as
    defined in Section~\ref{sec:ft-notification:coll}), it can
    continue to raise an error even after the failure has been
    acknowledged.
\end{users}

\section{Fault Tolerance Error Codes and Classes}
\label{sec:ft-errorcodes}
\mpitermindex{error handling!fault tolerance!notification}
\mpitermtitleindex{error handling!fault tolerance!fault tolerance error}

Among the error classes defined in Section~\ref{sec:ei-error-classes}, the following are \mpitermdefni{fault tolerance error} classes:
\begin{table}[htb]
\begin{center}
\begin{tabular}{l p{2.8in}}
    \error{MPI\_ERR\_PROC\_FAILED} & The operation could not complete because of an \MPI/ process failure (a fail-stop failure). \\
    \error{MPI\_ERR\_PROC\_FAILED\_PENDING} & The operation was interrupted by an \MPI/ process failure (a fail-stop failure). The request is still pending and the operation may be completed later. \\
    \error{MPI\_ERR\_REVOKED} & The communication object used in the operation has been revoked. \\
\end{tabular}
\end{center}
\caption{Fault tolerance error classes}
\label{table:ft-errorcodes:errclasses}
\end{table}

\section{Examples}

\subsection{Fault Tolerant Manager/Worker}

The example below presents a manager code that handles worker failures
by discarding failed worker \MPI/ processes and resubmitting the work to
the remaining workers. It demonstrates the different failure cases
that may occur with receive operations using \mpiconst{MPI\_ANY\_SOURCE}
as discussed in Section~\ref{sec:ft-notification:p2p-comm}.

\begin{example}
\label{ft-ex-manager-worker}
\Exindex{C}{Fault Tolerant Manager/Worker}{MPI\_Comm\_get\_failed,MPI\_Comm\_ack\_failed}%
Fault Tolerant Manager Example
%%HEADER
%%LANG: C
%%FRAGMENT
%%AUTOCONTINUE
%%DECL: MPI_Comm comm;
%%DECL: MPI_Group gcomm, g;
%%DECL: int rc, ec, size, gsize, tag, active_workers, work_available;
%%DECL: int *granks, *cranks;
%%DECL: MPI_Status status;
%%DECL: MPI_Request req;
%%ENDHEADER
\begin{lstlisting}[language={[MPI]C}]]
int manager(void)
{
    MPI_Comm_set_errhandler(comm, MPI_ERRORS_RETURN);
    MPI_Comm_size(comm, &size);
    MPI_Comm_group(comm, &gcomm);

    /* ... submit the initial work requests ... */

    /* Progress engine: Get answers, send new requests,
       and handle process failures */
    MPI_Irecv(buffer, 1, MPI_INT, MPI_ANY_SOURCE, tag, comm, &req);
    while( (active_workers > 0) && work_available ) {
        rc = MPI_Wait(&req, &status);
        if( MPI_SUCCESS == rc ) {
            /* ... process the answer and update work_available ... */
        }
        else {
            MPI_Error_class(rc, &ec);
            if( (MPI_ERR_PROC_FAILED == ec) ||
                (MPI_ERR_PROC_FAILED_PENDING == ec) ) {
                /* We ack the full size of comm, so we will ack
                 * unconditionally. Variable gsize will contain all
                 * currently known failures. */
                MPI_Comm_ack_failed(comm, size, &gsize);

                /* ... find the lost work and requeue it ... */
                MPI_Comm_get_failed(comm, &g);
                granks = (int*)calloc(active_workers-gsize-1,
                                      sizeof(int));
                cranks = (int*)calloc(active_workers-gsize-1,
                                      sizeof(int));
                for(i = active_workers; i < gsize; i++)
                    granks[i-active_workers] = i;
                MPI_Group_translate_ranks(g, gsize, granks,
                                          gcomm,    cranks);
                /* iterate over newly failed procs */
                for(i = active_workers; i < gsize; i++) {
                    /* resubmit the work */
                }
                free(cranks); free(granks);
                MPI_Group_free(&g);

                active_workers = size - gsize - 1;

                /* no need to repost: the request is still pending */
                if( ec == MPI_ERR_PROC_FAILED_PENDING )
                    continue;
            }
        }
        /* get ready to receive more notifications from workers */
        MPI_Irecv(buffer, 1, MPI_INT, MPI_ANY_SOURCE, tag, comm, &req);
    }
    /* ... cancel request and cleanup ... */
}
\end{lstlisting}
\end{example}


